% !TEX root = ../7.RETENTIONTIME.tex
\section{DISCUSSION}

\subsection{Compare the retention time distribution function of theory and experiment}

\begin{itemize}
    \item \textbf{1 tank system}: As seen in Figure 3, the retention time distribution function of the experiment is totally different compared with the theory. Generally, the $E$ values of the theory are smaller than the experiment. Besides, the compact time values ($\theta$) of theory are larger than the experiment ones. This suggests that according to the theory, dead zones or bypasses might appear. However, since the tank has a simple design, bypasses are unlikely; the discrepancy is likely due to dead zones where the agitators cannot reach.
    
    \item \textbf{2 tank system}: Following Figure 4, the retention time distribution function of the experiment somewhat resembles the theory (the $E$ value increases to a maximum and then decreases). However, the increase is not regular and is quite fluctuated. Additionally, the $E$ values of the theory are higher than the experiment.
    
    \item \textbf{3 tank system}: The retention time distribution function of the experiment is completely different compared with the theory and is extremely fluctuated.
\end{itemize}

\subsection{Phenomena causing unsteady state in process and equipment}
\begin{itemize}
    \item \textbf{Dead zones}: Created because the speed of the agitators is too slow, leaving some volumes unstirred.
    \item \textbf{Fluid circulation}: Caused by uneven mixing due to fluid properties or slow agitator speed.
    \item \textbf{Overcoming phenomena}: These can be mitigated by increasing the agitator speed for more uniform mixing.
    \item \textbf{Other factors}:
    \begin{itemize}
        \item Alteration of speed and movement direction due to equipment geometry (reactor, agitators, piping), creating swirl or dead areas.
        \item Unstable rotor speed over time (due to motor or voltage fluctuations).
        \item Indicator absorption onto the equipment walls.
        \item Ambient temperature fluctuations altering fluid properties (viscosity, density, velocity).
    \end{itemize}
\end{itemize}

\subsection{Evaluation of Errors}
\subsubsection{Error due to equipment}
\begin{itemize}
    \item \textbf{Photometer}: This is the primary source of error.
    \begin{itemize}
        \item Dirty cuvettes or samples absorb light, influencing results.
        \item The photometer is highly sensitive; even minor smudges on the cuvette surface affect readings.
        \item Aging of the machine and unstable power supply.
    \end{itemize}
    \item \textbf{Agitators}: Uneven swirling and unsteady velocity due to voltage fluctuations.
    \item \textbf{Piping}: Resistance in valves and piping makes it difficult to maintain constant flowrates, especially in multi-tank systems. Volume instability in the tanks during the experiment affects the steady state.
\end{itemize}

\subsubsection{Operation errors}
\begin{itemize}
    \item Timing precision when taking samples.
    \item Improper cleaning or drying of cuvettes, leading to contamination between samples.
    \item Unstable control of flowrates.
\end{itemize}

\subsubsection{Errors because of initial moment and pulse signal}
It is difficult to establish the exact Dirac delta function because the injection of the indicator takes a finite amount of time (approx. 5 seconds), which is significant and causes discrepancies between the theoretical and experimental curves.
